\documentclass[12pt]{article}
\usepackage{amsmath, amssymb}
\usepackage[margin=1in]{geometry}
\setlength{\parskip}{6pt}
\title{QUBO Formulation: Brain-Inspired Binary Sparse Coding Problem}
\date{}
\begin{document}
\maketitle

\subsection*{Brain-Inspired Binary Sparse Coding Problem}

\paragraph{Problem Statement.}
Given an input pattern vector $y \in \mathbb{R}^M$ and an overcomplete dictionary of basis patterns $D = \{d_0, d_1, \dots, d_{N-1}\}$ where each $d_i \in \mathbb{R}^M$, the goal is to select a sparse binary subset of dictionary atoms that reconstructs the input pattern with minimal error. The reconstruction is formed by a linear combination of the selected atoms. The objective is to minimize the sum of the squared reconstruction error and a sparsity penalty proportional to the number of selected atoms.

\paragraph{Decision Variables.}
For each dictionary atom $i \in \{0, 1, \dots, N-1\}$, we introduce a binary decision variable $x_i \in \{0, 1\}$. The variable $x_i = 1$ indicates that the $i$-th dictionary atom $d_i$ is selected for the reconstruction, and $x_i = 0$ indicates it is not selected.

\paragraph{QUBO Derivation.}
We derive the Quadratic Unconstrained Binary Optimization (QUBO) formulation by expanding the objective function and collecting terms.

\textbf{Step 1: Original Objective.}
The objective is to minimize the reconstruction error plus a sparsity penalty:
\begin{align}
    \min_{x} \left\| y - \sum_{i=0}^{N-1} x_i d_i \right\|^2 + \lambda \sum_{i=0}^{N-1} x_i,
\end{align}
where $\lambda > 0$ is the sparsity regularization parameter. Expanding the squared norm yields:
\begin{align}
    \left\| y - \sum_{i=0}^{N-1} x_i d_i \right\|^2 &= \left( y - \sum_{i=0}^{N-1} x_i d_i \right)^T \left( y - \sum_{j=0}^{N-1} x_j d_j \right) \\
    &= y^T y - 2 y^T \sum_{i=0}^{N-1} x_i d_i + \left( \sum_{i=0}^{N-1} x_i d_i \right)^T \left( \sum_{j=0}^{N-1} x_j d_j \right) \\
    &= \|y\|^2 - 2 \sum_{i=0}^{N-1} (y^T d_i) x_i + \sum_{i=0}^{N-1} \sum_{j=0}^{N-1} (d_i^T d_j) x_i x_j.
\end{align}
Since $\|y\|^2$ is constant with respect to $x$, it can be omitted from the optimization. The objective becomes:
\begin{align}
    H(x) = \sum_{i=0}^{N-1} \sum_{j=0}^{N-1} (d_i^T d_j) x_i x_j - 2 \sum_{i=0}^{N-1} (y^T d_i) x_i + \lambda \sum_{i=0}^{N-1} x_i.
\end{align}

\textbf{Step 2: Constraints.}
The problem is unconstrained; there are no constraints on the selection of atoms. Thus, no penalty terms are required.

\textbf{Step 3: Penalty Terms.}
As there are no constraints, the penalty term is zero.

\textbf{Step 4: Combined QUBO Hamiltonian.}
We combine the terms from Step 1. Using the idempotent property $x_i^2 = x_i$ for binary variables, we separate the quadratic and linear contributions. The double sum can be split into diagonal ($i=j$) and off-diagonal ($i \neq j$) terms:
\begin{align}
    H(x) &= \sum_{i=0}^{N-1} (d_i^T d_i) x_i^2 + \sum_{i=0}^{N-1} \sum_{j \neq i} (d_i^T d_j) x_i x_j - 2 \sum_{i=0}^{N-1} (y^T d_i) x_i + \lambda \sum_{i=0}^{N-1} x_i \\
    &= \sum_{i=0}^{N-1} (d_i^T d_i) x_i + \sum_{i=0}^{N-1} \sum_{j \neq i} (d_i^T d_j) x_i x_j - 2 \sum_{i=0}^{N-1} (y^T d_i) x_i + \lambda \sum_{i=0}^{N-1} x_i.
\end{align}
Grouping the linear terms for each $x_i$:
\begin{align}
    H(x) &= \sum_{i=0}^{N-1} \left[ (d_i^T d_i) - 2 (y^T d_i) + \lambda \right] x_i + \sum_{i=0}^{N-1} \sum_{j \neq i} (d_i^T d_j) x_i x_j.
\end{align}
To express this in standard QUBO form $H(x) = \sum_{i \leq j} Q_{i,j} x_i x_j + c$, we note that the off-diagonal sum counts each pair twice. We rewrite the off-diagonal sum as $2 \sum_{i < j} (d_i^T d_j) x_i x_j$. Thus:
\begin{align}
    H(x) = \sum_{i=0}^{N-1} \left[ (d_i^T d_i) - 2 (y^T d_i) + \lambda \right] x_i + \sum_{0 \leq i < j \leq N-1} 2 (d_i^T d_j) x_i x_j.
\end{align}

\textbf{Step 5: Q Matrix Entries.}
Reading off the coefficients from the expression above, the QUBO matrix entries $Q_{i,j}$ and offset $c$ are given by:
\begin{align}
    Q_{i,i} &= (d_i^T d_i) - 2 (y^T d_i) + \lambda, \quad \text{for } i \in \{0, \dots, N-1\}, \\
    Q_{i,j} &= 2 (d_i^T d_j), \quad \text{for } 0 \leq i < j \leq N-1, \\
    c &= 0.
\end{align}

\paragraph{Penalty Weight Justification.}
The formulation is unconstrained; therefore, no penalty weights are required to enforce constraints. The parameter $\lambda$ serves as the sparsity regularization strength. It is chosen based on the desired trade-off between reconstruction fidelity and sparsity. A larger $\lambda$ penalizes the selection of atoms more heavily, encouraging a sparser solution. In practice, $\lambda$ is often tuned via cross-validation or set proportional to the noise level in the input pattern $y$. Since there are no constraints, $\lambda$ does not need to satisfy a dominance condition over constraint violation penalties.

\paragraph{Correctness Argument.}
Minimizing the Hamiltonian $H(x)$ over the domain $\{0,1\}^N$ recovers the optimal sparse coding solution because the QUBO formulation is an exact algebraic transformation of the original objective function. The quadratic terms $\sum_{i<j} 2(d_i^T d_j)x_i x_j$ correctly capture the interference between pairs of selected atoms in the reconstruction error, while the linear terms $(d_i^T d_i) - 2(y^T d_i) + \lambda$ account for the self-interaction of each atom, its alignment with the input, and the sparsity cost. Since the mapping is exact and unconstrained, the global minimum of $H(x)$ corresponds precisely to the binary vector $x$ that minimizes the original reconstruction error plus sparsity penalty.

\paragraph{Illustrative Example.}
Consider a small instance with $N=3$ dictionary atoms. The variables are $x_0, x_1, x_2 \in \{0,1\}$. The QUBO matrix entries are instantiated from the general formulas as follows:
\begin{align}
    Q_{0,0} &= (d_0^T d_0) - 2 (y^T d_0) + \lambda, \\
    Q_{1,1} &= (d_1^T d_1) - 2 (y^T d_1) + \lambda, \\
    Q_{2,2} &= (d_2^T d_2) - 2 (y^T d_2) + \lambda, \\
    Q_{0,1} &= 2 (d_0^T d_1), \\
    Q_{0,2} &= 2 (d_0^T d_2), \\
    Q_{1,2} &= 2 (d_1^T d_2).
\end{align}
The resulting Hamiltonian is:
\begin{align}
    H(x) &= Q_{0,0} x_0 + Q_{1,1} x_1 + Q_{2,2} x_2 + 2 Q_{0,1} x_0 x_1 + 2 Q_{0,2} x_0 x_2 + 2 Q_{1,2} x_1 x_2.
\end{align}
Substituting the expressions for $Q_{i,j}$:
\begin{align}
    H(x) &= \left[ (d_0^T d_0) - 2 (y^T d_0) + \lambda \right] x_0 + \left[ (d_1^T d_1) - 2 (y^T d_1) + \lambda \right] x_1 + \left[ (d_2^T d_2) - 2 (y^T d_2) + \lambda \right] x_2 \nonumber \\
    &\quad + 2 (d_0^T d_1) x_0 x_1 + 2 (d_0^T d_2) x_0 x_2 + 2 (d_1^T d_2) x_1 x_2.
\end{align}
The offset is $c=0$. The optimal solution $x^*$ is the binary vector that minimizes this energy, yielding the selected dictionary atoms and the minimum sparse-coding energy.

\end{document}